\input includes/\format

\pdfinfo{
   /Title (Gokyo Seimyo)
   /Creator (TeX)
   /Author (Vit Brunner)
   /CreationDate (20261111233500)
}

~
\title
\parindent=0em
\baselineskip=1em

\vskip 5cm
\centerline{\japanesetitle 碁経精妙}
\bigskip
\centerline{Gokyo Seimyo}
\vfill\break

\input includes/\format-motto

\bigskip
{\header Preface}
\medskip
Gokyo Seimyo was published by Hayashi Genbi in 1835 as a sequel to Gokyo Shumyo. This booklet contains 145 diagrams from Volume 3, three of which are reference diagrams, and the rest are problems to solve. Problems are shown as {\bf black to play}.
\bigskip

\def\chapterentry#1#2#3{%
  \pdfsectionbookmark{#1}{#2}%
  \hbox to \hsize{\strut
    \pdfstartlink attr{/Border [0 0 0]} goto name {problem-#1-1}%
    #1. #2 (#3 diagrams)\dotfill \chapterpage{#1}%
    \pdfendlink}}

\vbox{
\noindent{\header Sections}
\medskip
\chapterentry{1}{Tesuji inside territories}{44}
\chapterentry{2}{Other problems}{74}
\chapterentry{3}{Connect underneath}{10}
\chapterentry{4}{Cut/connect groups}{8}
\chapterentry{5}{Ladder}{9}

\bigskip
\rightline{\it\me}
\rightline{\it November 2026}
}
\vfill\break

\input includes/\format-gotc
\baselineskip=1.1em
\input generated/\format-gokyo-seimyo

\bye
